\documentclass[12pt]{article}

\usepackage[a4paper,margin=1in]{geometry}
\usepackage{amsmath,amssymb}
\usepackage{newtxtext,newtxmath}
\usepackage{fancyhdr}
\usepackage{listings}
\usepackage{xcolor}
\usepackage{graphicx}
\usepackage{float}
\usepackage{booktabs}
\usepackage[T1]{fontenc}

\setlength{\textfloatsep}{10pt}
\setlength{\floatsep}{10pt}
\setlength{\parindent}{0pt}
\setlength{\parskip}{5pt}
\AtBeginDocument{\fontsize{12pt}{14pt}\selectfont}

\lstdefinestyle{code}{
    language=C++,
    basicstyle=\ttfamily\fontsize{7.2pt}{8.5pt}\selectfont,
    breaklines=true,
    breakatwhitespace=false,
    frame=single,
    numbers=left,
    numberstyle=\tiny\color{gray},
    keywordstyle=\color{blue},
    commentstyle=\color{gray},
    stringstyle=\color{green!50!black},
    columns=fullflexible,
    keepspaces=true,
    showstringspaces=false,
    tabsize=4
}
\lstdefinestyle{output}{
    basicstyle=\ttfamily\fontsize{8.2pt}{9.5pt}\selectfont,
    breaklines=true,
    frame=single,
    numbers=none,
    backgroundcolor=\color{gray!6},
    columns=fullflexible,
    keepspaces=true,
    showstringspaces=false
}

\pagestyle{fancy}
\fancyhf{}
\fancyhead[L]{Name: Kishan Sahu}
\fancyhead[R]{Enrollment No.: P26DS017}
\renewcommand{\headrulewidth}{0.4pt}
\setlength{\headheight}{15pt}

\begin{document}

\begin{center}
\textbf{Computer Science and Engineering Department, SVNIT, Surat}\\
\textbf{M.Tech. I -- Semester I}\\
\textbf{High Performance Computing (CSDS125)}\\
\textbf{Lab Assignment 8: Parallel Matrix Multiplication and Sobel Edge Detection using OpenMPI}
\end{center}
\vspace{-4pt}
\hrule
\vspace{6pt}

\section*{Part 1: Parallel Tile-Based Matrix Multiplication (Cannon's Algorithm)}

\subsection*{1. Problem Statement and Algorithm Overview}
The objective is to implement parallel matrix multiplication of two $N \times N$ matrices, $A$ and $B$, to compute $C = A \times B$ using OpenMPI in C. To achieve scalable load balance and minimize communication overhead, Cannon's algorithm organizes $P = q^2$ processes into a $q \times q$ 2D Cartesian grid:
\begin{enumerate}
    \item \textbf{Topology Setup:} A 2D periodic Cartesian communicator is constructed using \texttt{MPI\_Cart\_create}. Each process $(i, j)$ is assigned square sub-matrix tiles $A_{i,j}$, $B_{i,j}$, and $C_{i,j}$ of size $b \times b$ where $b = N / q$.
    \item \textbf{Initial Skewing (Alignment):} Process $(i, j)$ circularly shifts tile $A_{i,j}$ to the left by $i$ positions (row shift) and tile $B_{i,j}$ upwards by $j$ positions (column shift) using \texttt{MPI\_Cart\_shift} and \texttt{MPI\_Sendrecv\_replace}.
    \item \textbf{Multiply-and-Shift Loop:} Over $q$ iterations:
    \begin{itemize}
        \item Each process performs local tile multiplication: $C_{i,j} \leftarrow C_{i,j} + A_{i,j} \times B_{i,j}$.
        \item Sub-matrix $A_{i,j}$ is circularly shifted left by 1 position, and $B_{i,j}$ is circularly shifted up by 1 position.
    \end{itemize}
    \item \textbf{Collection and Verification:} Rank 0 gathers the computed tiles, reconstructs the global matrix $C$, and verifies accuracy against a sequential matrix multiplication baseline within tolerance $\epsilon = 10^{-5}$.
\end{enumerate}

\subsection*{2. Complete Source Code: \texttt{tile\_matmul.c}}
\begin{lstlisting}[style=code]
#include <stdio.h>
#include <stdlib.h>
#include <math.h>
#include <mpi.h>

// Multiply two b x b tiles and accumulate into C_tile: C += A * B
void multiply_tile(int b, double *A, double *B, double *C) {
    for (int i = 0; i < b; i++) {
        for (int k = 0; k < b; k++) {
            for (int j = 0; j < b; j++) {
                C[i * b + j] += A[i * b + k] * B[k * b + j];
            }
        }
    }
}

// Print an N x N matrix
void print_matrix(const char *name, double *mat, int N) {
    printf("%s (%dx%d):\n", name, N, N);
    for (int i = 0; i < N; i++) {
        for (int j = 0; j < N; j++) {
            printf("%7.2f ", mat[i * N + j]);
        }
        printf("\n");
    }
    printf("\n");
}

int main(int argc, char *argv[]) {
    int rank, size;
    MPI_Init(&argc, &argv);
    MPI_Comm_size(MPI_COMM_WORLD, &size);
    MPI_Comm_rank(MPI_COMM_WORLD, &rank);

    // Matrix dimension N (default 4)
    int N = 4;
    if (argc > 1) {
        N = atoi(argv[1]);
    }

    // Check if process count P is a perfect square: P = q * q
    int q = (int)sqrt(size);
    if (q * q != size) {
        if (rank == 0) {
            fprintf(stderr, "Error: Number of processes (%d) must be a perfect square (e.g., 1, 4, 9, 16).\n", size);
        }
        MPI_Abort(MPI_COMM_WORLD, 1);
    }

    // Check if matrix dimension N is divisible by grid dimension q
    if (N % q != 0) {
        if (rank == 0) {
            fprintf(stderr, "Error: Matrix dimension N (%d) must be divisible by grid dimension q (%d).\n", N, q);
        }
        MPI_Abort(MPI_COMM_WORLD, 1);
    }

    int b = N / q; // tile size (b x b)
    int tile_size = b * b;

    // Create 2D Cartesian Grid Topology
    int dims[2] = {q, q};
    int periods[2] = {1, 1}; // periodic/circular boundaries in both dimensions
    MPI_Comm cart_comm;
    MPI_Cart_create(MPI_COMM_WORLD, 2, dims, periods, 1, &cart_comm);

    int coords[2];
    MPI_Cart_coords(cart_comm, rank, 2, coords);
    int row = coords[0];
    int col = coords[1];

    // Allocate memory for local tiles
    double *A_tile = (double *)malloc(tile_size * sizeof(double));
    double *B_tile = (double *)malloc(tile_size * sizeof(double));
    double *C_tile = (double *)calloc(tile_size, sizeof(double));

    // Global matrices on rank 0
    double *A_global = NULL;
    double *B_global = NULL;
    double *C_global = NULL;

    if (rank == 0) {
        A_global = (double *)malloc(N * N * sizeof(double));
        B_global = (double *)malloc(N * N * sizeof(double));
        C_global = (double *)malloc(N * N * sizeof(double));

        // Initialize sample matrices
        for (int i = 0; i < N; i++) {
            for (int j = 0; j < N; j++) {
                A_global[i * N + j] = i + j + 1.0;
                B_global[i * N + j] = (i == j) ? 2.0 : 1.0;
            }
        }

        if (N <= 8) {
            print_matrix("Matrix A", A_global, N);
            print_matrix("Matrix B", B_global, N);
        }

        // Distribute tiles to each process
        for (int p = 0; p < size; p++) {
            int p_coords[2];
            MPI_Cart_coords(cart_comm, p, 2, p_coords);
            int pr = p_coords[0];
            int pc = p_coords[1];

            double temp_A[tile_size];
            double temp_B[tile_size];

            for (int i = 0; i < b; i++) {
                for (int j = 0; j < b; j++) {
                    int g_idx = (pr * b + i) * N + (pc * b + j);
                    temp_A[i * b + j] = A_global[g_idx];
                    temp_B[i * b + j] = B_global[g_idx];
                }
            }

            if (p == 0) {
                for (int k = 0; k < tile_size; k++) {
                    A_tile[k] = temp_A[k];
                    B_tile[k] = temp_B[k];
                }
            } else {
                MPI_Send(temp_A, tile_size, MPI_DOUBLE, p, 0, cart_comm);
                MPI_Send(temp_B, tile_size, MPI_DOUBLE, p, 1, cart_comm);
            }
        }
    } else {
        MPI_Recv(A_tile, tile_size, MPI_DOUBLE, 0, 0, cart_comm, MPI_STATUS_IGNORE);
        MPI_Recv(B_tile, tile_size, MPI_DOUBLE, 0, 1, cart_comm, MPI_STATUS_IGNORE);
    }

    // Barrier and Start Timer
    MPI_Barrier(cart_comm);
    double start_time = MPI_Wtime();

    // Cannon's Algorithm - Phase 1: Initial Alignment (Skewing)
    if (row > 0) {
        int src_A, dst_A;
        MPI_Cart_shift(cart_comm, 1, -row, &src_A, &dst_A);
        MPI_Sendrecv_replace(A_tile, tile_size, MPI_DOUBLE, dst_A, 10, src_A, 10, cart_comm, MPI_STATUS_IGNORE);
    }

    if (col > 0) {
        int src_B, dst_B;
        MPI_Cart_shift(cart_comm, 0, -col, &src_B, &dst_B);
        MPI_Sendrecv_replace(B_tile, tile_size, MPI_DOUBLE, dst_B, 20, src_B, 20, cart_comm, MPI_STATUS_IGNORE);
    }

    // Cannon's Algorithm - Phase 2: Compute and Shift (q steps)
    int src_A_step, dst_A_step;
    int src_B_step, dst_B_step;
    MPI_Cart_shift(cart_comm, 1, -1, &src_A_step, &dst_A_step); // left shift by 1
    MPI_Cart_shift(cart_comm, 0, -1, &src_B_step, &dst_B_step); // up shift by 1

    for (int step = 0; step < q; step++) {
        // Multiply current local tiles
        multiply_tile(b, A_tile, B_tile, C_tile);

        // Circularly shift A left by 1 and B up by 1 (except after last step)
        if (step < q - 1) {
            MPI_Sendrecv_replace(A_tile, tile_size, MPI_DOUBLE, dst_A_step, 30, src_A_step, 30, cart_comm, MPI_STATUS_IGNORE);
            MPI_Sendrecv_replace(B_tile, tile_size, MPI_DOUBLE, dst_B_step, 40, src_B_step, 40, cart_comm, MPI_STATUS_IGNORE);
        }
    }

    // Barrier and Stop Timer
    MPI_Barrier(cart_comm);
    double end_time = MPI_Wtime();
    double elapsed_time = end_time - start_time;

    // Collect result tiles on rank 0
    if (rank == 0) {
        for (int p = 0; p < size; p++) {
            double temp_C[tile_size];
            if (p == 0) {
                for (int k = 0; k < tile_size; k++) {
                    temp_C[k] = C_tile[k];
                }
            } else {
                MPI_Recv(temp_C, tile_size, MPI_DOUBLE, p, 50, cart_comm, MPI_STATUS_IGNORE);
            }

            int p_coords[2];
            MPI_Cart_coords(cart_comm, p, 2, p_coords);
            int pr = p_coords[0];
            int pc = p_coords[1];

            for (int i = 0; i < b; i++) {
                for (int j = 0; j < b; j++) {
                    C_global[(pr * b + i) * N + (pc * b + j)] = temp_C[i * b + j];
                }
            }
        }
    } else {
        MPI_Send(C_tile, tile_size, MPI_DOUBLE, 0, 50, cart_comm);
    }

    // Verification and Report on rank 0
    if (rank == 0) {
        if (N <= 8) {
            print_matrix("Result Matrix C (Parallel)", C_global, N);
        }

        // Sequential reference computation
        double *C_seq = (double *)calloc(N * N, sizeof(double));
        for (int i = 0; i < N; i++) {
            for (int k = 0; k < N; k++) {
                for (int j = 0; j < N; j++) {
                    C_seq[i * N + j] += A_global[i * N + k] * B_global[k * N + j];
                }
            }
        }

        // Verify correctness
        int correct = 1;
        for (int i = 0; i < N * N; i++) {
            if (fabs(C_global[i] - C_seq[i]) > 1e-5) {
                correct = 0;
                break;
            }
        }

        printf("==================================================\n");
        printf("Parallel Tile Matrix Multiplication (Cannon's Algorithm)\n");
        printf("Processes: %d (Grid: %dx%d)\n", size, q, q);
        printf("Matrix Dimension: %dx%d (Tile Size: %dx%d)\n", N, N, b, b);
        printf("Elapsed Time: %f seconds\n", elapsed_time);
        printf("Verification: %s\n", correct ? "PASSED" : "FAILED");
        printf("==================================================\n");

        free(A_global);
        free(B_global);
        free(C_global);
        free(C_seq);
    }

    // Clean up local allocations
    free(A_tile);
    free(B_tile);
    free(C_tile);
    MPI_Comm_free(&cart_comm);

    MPI_Finalize();
    return 0;
}
\end{lstlisting}

\subsection*{3. Terminal Execution Output: Matrix Multiplication}
\begin{lstlisting}[style=output]
Matrix A (4x4):
   1.00    2.00    3.00    4.00 
   2.00    3.00    4.00    5.00 
   3.00    4.00    5.00    6.00 
   4.00    5.00    6.00    7.00 

Matrix B (4x4):
   2.00    1.00    1.00    1.00 
   1.00    2.00    1.00    1.00 
   1.00    1.00    2.00    1.00 
   1.00    1.00    1.00    2.00 

Result Matrix C (Parallel) (4x4):
  11.00   12.00   13.00   14.00 
  16.00   17.00   18.00   19.00 
  21.00   22.00   23.00   24.00 
  26.00   27.00   28.00   29.00 

==================================================
Parallel Tile Matrix Multiplication (Cannon's Algorithm)
Processes: 4 (Grid: 2x2)
Matrix Dimension: 4x4 (Tile Size: 2x2)
Elapsed Time: 0.004005 seconds
Verification: PASSED
==================================================
\end{lstlisting}

\vspace{10pt}
\hrule
\vspace{10pt}

\section*{Part 2: Parallel Sobel-Filter Edge Detection with Halo Regions}

\subsection*{1. Problem Statement and Algorithm Overview}
The Sobel operator is an edge detection algorithm that computes gradient magnitude approximations across an image using two $3 \times 3$ convolution kernels:
\begin{equation}
G_x = \begin{bmatrix}
-1 & 0 & 1 \\
-2 & 0 & 2 \\
-1 & 0 & 1
\end{bmatrix}, \quad
G_y = \begin{bmatrix}
-1 & -2 & -1 \\
0 & 0 & 0 \\
1 & 2 & 1
\end{bmatrix}
\end{equation}
The gradient magnitude at each pixel $(i, j)$ is given by $G = \min(255, \operatorname{round}(\sqrt{G_x^2 + G_y^2}))$.

\textbf{Concept of Halo Region (Ghost Cells):}
When an $H \times W$ image is partitioned horizontally among $P$ processes, each process holds $H_{local} = H / P$ rows. Computing the $3 \times 3$ stencil on row 0 of a local chunk requires the row above it, which belongs to rank $r - 1$. Similarly, computing on row $H_{local} - 1$ requires the row below, belonging to rank $r + 1$.
\begin{itemize}
    \item \textbf{Halo Buffer Allocation:} Each rank allocates an extended local buffer of size $(H_{local} + 2) \times W$. Row $0$ acts as the top halo row, rows $1$ through $H_{local}$ store active data, and row $H_{local} + 1$ acts as the bottom halo row.
    \item \textbf{Deadlock-Free Halo Exchange:} Adjacent ranks exchange border rows using \texttt{MPI\_Sendrecv}. Rank 0 uses \texttt{MPI\_PROC\_NULL} for its top neighbor, and rank $P - 1$ uses \texttt{MPI\_PROC\_NULL} for its bottom neighbor (with zero-padding applied to global boundaries).
    \item \textbf{Finer Output:} Utilizing halo rows eliminates boundary distortion and row loss across partition lines, yielding 100\% mathematical parity with sequential full-image filtering.
\end{itemize}

\subsection*{2. Complete Source Code: \texttt{sobel\_edge.c}}
\begin{lstlisting}[style=code]
#include <stdio.h>
#include <stdlib.h>
#include <string.h>
#include <ctype.h>
#include <math.h>
#include <mpi.h>

// Save grayscale image to a standard binary PGM (P5) file
void save_pgm(const char *filename, unsigned char *img, int H, int W) {
    FILE *f = fopen(filename, "wb");
    if (!f) {
        fprintf(stderr, "Warning: Could not open %s for writing.\n", filename);
        return;
    }
    fprintf(f, "P5\n%d %d\n255\n", W, H);
    fwrite(img, sizeof(unsigned char), H * W, f);
    fclose(f);
}

// Helper to skip whitespace and comments in PGM header
static void skip_pgm_comments(FILE *f) {
    int ch;
    while ((ch = fgetc(f)) != EOF) {
        if (ch == '#') {
            while ((ch = fgetc(f)) != EOF && ch != '\n');
        } else if (!isspace(ch)) {
            ungetc(ch, f);
            break;
        }
    }
}

// Load a binary PGM (P5) image from disk
unsigned char *load_pgm(const char *filename, int *H, int *W) {
    FILE *f = fopen(filename, "rb");
    if (!f) {
        return NULL;
    }

    char magic[3];
    if (fscanf(f, "%2s", magic) != 1 || strcmp(magic, "P5") != 0) {
        fclose(f);
        return NULL;
    }

    skip_pgm_comments(f);
    if (fscanf(f, "%d", W) != 1) {
        fclose(f);
        return NULL;
    }

    skip_pgm_comments(f);
    if (fscanf(f, "%d", H) != 1) {
        fclose(f);
        return NULL;
    }

    int max_val;
    skip_pgm_comments(f);
    if (fscanf(f, "%d", &max_val) != 1 || max_val > 255) {
        fclose(f);
        return NULL;
    }

    fgetc(f); // skip single whitespace

    unsigned char *data = (unsigned char *)malloc((*H) * (*W) * sizeof(unsigned char));
    if (!data) {
        fclose(f);
        return NULL;
    }

    size_t read_bytes = fread(data, sizeof(unsigned char), (*H) * (*W), f);
    if (read_bytes != (size_t)((*H) * (*W))) {
        fprintf(stderr, "Warning: Expected %d bytes, read %zu bytes from %s.\n", (*H) * (*W), read_bytes, filename);
    }

    fclose(f);
    return data;
}

// Generate a synthetic test image with geometric features and sharp edges
void generate_synthetic_image(unsigned char *img, int H, int W) {
    for (int i = 0; i < H * W; i++) {
        img[i] = 30;
    }

    int r_start = H / 4, r_end = 3 * H / 4;
    int c_start = W / 4, c_end = 3 * W / 4;
    for (int i = r_start; i < r_end; i++) {
        for (int j = c_start; j < c_end; j++) {
            img[i * W + j] = 220;
        }
    }

    int cy = H / 2, cx = W / 2;
    int radius = H / 6;
    for (int i = 0; i < H; i++) {
        for (int j = 0; j < W; j++) {
            int dy = i - cy;
            int dx = j - cx;
            if (dy * dy + dx * dx <= radius * radius) {
                img[i * W + j] = 90;
            }
        }
    }
}

// Sequential Sobel reference computation on full image (for verification)
void sequential_sobel(unsigned char *src, unsigned char *dst, int H, int W) {
    for (int i = 0; i < H; i++) {
        for (int j = 0; j < W; j++) {
            if (i == 0 || i == H - 1 || j == 0 || j == W - 1) {
                dst[i * W + j] = 0;
            } else {
                int gx = -1 * src[(i - 1) * W + (j - 1)] + 1 * src[(i - 1) * W + (j + 1)]
                         -2 * src[i * W + (j - 1)]       + 2 * src[i * W + (j + 1)]
                         -1 * src[(i + 1) * W + (j - 1)] + 1 * src[(i + 1) * W + (j + 1)];

                int gy = -1 * src[(i - 1) * W + (j - 1)] - 2 * src[(i - 1) * W + j] - 1 * src[(i - 1) * W + (j + 1)]
                         +1 * src[(i + 1) * W + (j - 1)] + 2 * src[(i + 1) * W + j] + 1 * src[(i + 1) * W + (j + 1)];

                int mag = (int)(sqrt(gx * gx + gy * gy) + 0.5);
                if (mag > 255) mag = 255;
                dst[i * W + j] = (unsigned char)mag;
            }
        }
    }
}

int main(int argc, char *argv[]) {
    int rank, size;
    MPI_Init(&argc, &argv);
    MPI_Comm_rank(MPI_COMM_WORLD, &rank);
    MPI_Comm_size(MPI_COMM_WORLD, &size);

    int H = 0, W = 0;
    unsigned char *full_image = NULL;
    char image_source[256] = "";

    if (rank == 0) {
        const char *filename = (argc > 1 && argv[1][0] != '-') ? argv[1] : "lena.pgm";
        full_image = load_pgm(filename, &H, &W);

        if (full_image) {
            snprintf(image_source, sizeof(image_source), "PGM file '%s'", filename);
        } else {
            H = (argc >= 2 && atoi(argv[1]) > 0) ? atoi(argv[1]) : 512;
            W = (argc >= 3 && atoi(argv[2]) > 0) ? atoi(argv[2]) : H;
            full_image = (unsigned char *)malloc(H * W * sizeof(unsigned char));
            generate_synthetic_image(full_image, H, W);
            snprintf(image_source, sizeof(image_source), "Synthetic test pattern (%dx%d)", W, H);
        }
    }

    // Broadcast dimensions (H, W) to all processes
    int dims[2];
    if (rank == 0) {
        dims[0] = H;
        dims[1] = W;
    }
    MPI_Bcast(dims, 2, MPI_INT, 0, MPI_COMM_WORLD);
    H = dims[0];
    W = dims[1];

    if (H % size != 0) {
        if (rank == 0) {
            fprintf(stderr, "Error: Image height (%d) must be divisible by process count (%d).\n", H, size);
            if (full_image) free(full_image);
        }
        MPI_Abort(MPI_COMM_WORLD, 1);
    }

    int H_local = H / size;

    if (rank == 0) {
        printf("=====================================================\n");
        printf(" Parallel Sobel Edge Detection (OpenMPI with Halo)  \n");
        printf("=====================================================\n");
        printf("Input source     : %s\n", image_source);
        printf("Image dimensions : %d x %d\n", W, H);
        printf("MPI processes    : %d\n", size);
        printf("Rows per process : %d\n", H_local);
        printf("-----------------------------------------------------\n");
    }

    unsigned char *parallel_output = NULL;
    unsigned char *seq_output = NULL;
    if (rank == 0) {
        parallel_output = (unsigned char *)malloc(H * W * sizeof(unsigned char));
        seq_output = (unsigned char *)malloc(H * W * sizeof(unsigned char));
    }

    unsigned char *local_buf = (unsigned char *)malloc((H_local + 2) * W * sizeof(unsigned char));
    unsigned char *local_out = (unsigned char *)malloc(H_local * W * sizeof(unsigned char));

    MPI_Barrier(MPI_COMM_WORLD);
    double t_start = MPI_Wtime();

    // 1. Scatter image chunks to rows 1..H_local of local_buf
    MPI_Scatter(full_image, H_local * W, MPI_UNSIGNED_CHAR,
                &local_buf[1 * W], H_local * W, MPI_UNSIGNED_CHAR,
                0, MPI_COMM_WORLD);

    // 2. Halo Exchange with top and bottom neighbors
    int top_neighbor = (rank == 0) ? MPI_PROC_NULL : rank - 1;
    int bottom_neighbor = (rank == size - 1) ? MPI_PROC_NULL : rank + 1;

    MPI_Sendrecv(&local_buf[1 * W], W, MPI_UNSIGNED_CHAR, top_neighbor, 0,
                 &local_buf[(H_local + 1) * W], W, MPI_UNSIGNED_CHAR, bottom_neighbor, 0,
                 MPI_COMM_WORLD, MPI_STATUS_IGNORE);

    MPI_Sendrecv(&local_buf[H_local * W], W, MPI_UNSIGNED_CHAR, bottom_neighbor, 1,
                 &local_buf[0 * W], W, MPI_UNSIGNED_CHAR, top_neighbor, 1,
                 MPI_COMM_WORLD, MPI_STATUS_IGNORE);

    if (rank == 0) {
        memset(&local_buf[0 * W], 0, W);
    }
    if (rank == size - 1) {
        memset(&local_buf[(H_local + 1) * W], 0, W);
    }

    // 3. Compute Sobel filter using the halo region
    for (int i = 1; i <= H_local; i++) {
        int global_row = rank * H_local + (i - 1);

        for (int j = 0; j < W; j++) {
            if (global_row == 0 || global_row == H - 1 || j == 0 || j == W - 1) {
                local_out[(i - 1) * W + j] = 0;
            } else {
                int gx = -1 * local_buf[(i - 1) * W + (j - 1)] + 1 * local_buf[(i - 1) * W + (j + 1)]
                         -2 * local_buf[i * W + (j - 1)]       + 2 * local_buf[i * W + (j + 1)]
                         -1 * local_buf[(i + 1) * W + (j - 1)] + 1 * local_buf[(i + 1) * W + (j + 1)];

                int gy = -1 * local_buf[(i - 1) * W + (j - 1)] - 2 * local_buf[(i - 1) * W + j] - 1 * local_buf[(i - 1) * W + (j + 1)]
                         +1 * local_buf[(i + 1) * W + (j - 1)] + 2 * local_buf[(i + 1) * W + j] + 1 * local_buf[(i + 1) * W + (j + 1)];

                int mag = (int)(sqrt(gx * gx + gy * gy) + 0.5);
                if (mag > 255) mag = 255;
                local_out[(i - 1) * W + j] = (unsigned char)mag;
            }
        }
    }

    // 4. Gather computed edge rows back to master process
    MPI_Gather(local_out, H_local * W, MPI_UNSIGNED_CHAR,
               parallel_output, H_local * W, MPI_UNSIGNED_CHAR,
               0, MPI_COMM_WORLD);

    MPI_Barrier(MPI_COMM_WORLD);
    double t_end = MPI_Wtime();

    // 5. Master verification and reporting
    if (rank == 0) {
        double elapsed_time = t_end - t_start;
        printf("Parallel execution time   : %.6f seconds\n", elapsed_time);

        double seq_start = MPI_Wtime();
        sequential_sobel(full_image, seq_output, H, W);
        double seq_end = MPI_Wtime();
        printf("Sequential reference time : %.6f seconds\n", seq_end - seq_start);

        int mismatches = 0;
        for (int i = 0; i < H * W; i++) {
            if (parallel_output[i] != seq_output[i]) {
                mismatches++;
            }
        }

        if (mismatches == 0) {
            printf("\n[VERIFICATION PASSED] Parallel output with Halo Exchange matches sequential result 100%%!\n");
        } else {
            printf("\n[VERIFICATION FAILED] Found %d mismatched pixels out of %d!\n", mismatches, H * W);
        }

        save_pgm("output_edges.pgm", parallel_output, H, W);
        printf("Output saved to output_edges.pgm (PGM P5 format)\n");
        printf("=====================================================\n");

        free(full_image);
        free(parallel_output);
        free(seq_output);
    }

    free(local_buf);
    free(local_out);

    MPI_Finalize();
    return 0;
}
\end{lstlisting}

\subsection*{3. Terminal Execution Output: Sobel Edge Detection}
\begin{lstlisting}[style=output]
mpirun --oversubscribe -np 4 ./sobel_edge lena.pgm
=====================================================
 Parallel Sobel Edge Detection (OpenMPI with Halo)  
=====================================================
Input source     : PGM file 'lena.pgm'
Image dimensions : 512 x 512
MPI processes    : 4
Rows per process : 128
-----------------------------------------------------
Parallel execution time   : 0.002959 seconds
Sequential reference time : 0.008928 seconds

[VERIFICATION PASSED] Parallel output with Halo Exchange matches sequential result 100%!
Output saved to output_edges.pgm (PGM P5 format)
=====================================================
\end{lstlisting}

\subsection*{4. Visual Results Comparison}
The halo-based domain decomposition computes seamless edge transitions across the 4 process slab boundaries without any artifacting or loss of border rows:

\begin{figure}[H]
\centering
\begin{minipage}{0.48\textwidth}
    \centering
    \includegraphics[width=0.92\linewidth]{lena.png}
    \caption{Input Grayscale Image (\texttt{lena.png}, $512 \times 512$)}
\end{minipage}\hfill
\begin{minipage}{0.48\textwidth}
    \centering
    \includegraphics[width=0.92\linewidth]{output_edges.png}
    \caption{Sobel Edge Detection Output (\texttt{output\_edges.png})}
\end{minipage}
\end{figure}

\subsection*{5. Performance and Scaling Analysis}
\begin{table}[H]
\centering
\small
\begin{tabular}{@{}ccccc@{}}
\toprule
\textbf{MPI Processes ($P$)} & \textbf{Rows per Rank} & \textbf{Parallel Time (s)} & \textbf{Sequential Time (s)} & \textbf{Verification Status} \\ \midrule
1 & 512 & 0.011468 & 0.009205 & 100\% Match (PASSED) \\
2 & 256 & 0.003137 & 0.003148 & 100\% Match (PASSED) \\
4 & 128 & 0.002959 & 0.008928 & 100\% Match (PASSED) \\
8 & 64  & 0.004731 & 0.004086 & 100\% Match (PASSED) \\ \bottomrule
\end{tabular}
\caption{Execution time and verification results across varying OpenMPI process counts on $512 \times 512$ image}
\end{table}

\end{document}

